\documentclass[12pt]{article}

% ==== topological-phases-of-matter : canonical macros (v1, 2026-07-14) ====
% Requires amsmath, amssymb, amsfonts (mathfrak, mathbf, mathcal all standard).
\usepackage{amsmath,amssymb,amsfonts}

% --- condensation / probes ---
\newcommand{\cond}[1]{\underline{#1}}          % condensation X |-> \cond{X}
\newcommand{\cg}[1]{\underline{#1}}            % condensed group (alias, for readability)
\newcommand{\Cont}{\operatorname{Cont}}        % Cont(S,X)
\newcommand{\Shape}{\operatorname{Shape}}      % shape / homotopy type

% --- interaction space & locality ---
\newcommand{\Int}{\mathcal{I}}                 % \Int_{d,G} interaction space
\newcommand{\Ffun}{F}                          % F-function (interaction decay)
\newcommand{\BF}{\mathcal{B}_{F}}              % interaction Banach space for F

% --- moduli stacks ---
\newcommand{\Ham}{\mathfrak{Ham}}              % \Ham_{d,G}
\newcommand{\Gap}{\mathfrak{Gap}}              % \Gap_{d,G}
\newcommand{\Phase}{\mathfrak{Phase}}          % \Phase_{d,G}
\newcommand{\Phases}{\operatorname{Phases}}    % \Phases_{d,G} = pi_0 Shape Phase
\newcommand{\Weq}{\mathcal{W}}                 % class of equivalences to invert
\newcommand{\st}{\mathrm{st}}                  % stabilization superscript
\newcommand{\hfix}[1]{h#1}                     % homotopy-fixed-point superscript, e.g. \Phase_d^{\hfix{\cg{G}}}
\newcommand{\quotstack}[2]{[#1/#2]}            % action stack [Y/\cg{G}]

% --- gap ---
\newcommand{\gapof}[1]{\operatorname{gap}(#1)} % gap(H)
\newcommand{\gapbound}{\Delta}                  % uniform gap bound

% --- stacking / spectrum ---
% \boxtimes is built in (stacking product)
\newcommand{\IPcond}{\mathbf{IP}^{\mathrm{cond}}}  % invertible condensed phase spectrum

% --- transitions / discriminant / charge ---
\newcommand{\Sig}{\Sigma}                       % \Sig_f gapless discriminant
\newcommand{\rc}{\partial\nu}                   % relative transition charge in E^{q+1}(B,U_f)

% --- observables & solid K-theory ---
\newcommand{\Kop}{K_{\mathrm{op}}}             % operator / topological K-theory
\newcommand{\Kalg}{K_{\mathrm{alg}}}           % algebraic K-theory
\newcommand{\Solid}{\operatorname{Solid}}      % solidification functor

% --- disorder ---
\newcommand{\dis}{\Omega}                       % disorder hull \Omega = F^{Z^d}
% ==== end canonical macros ====

% ---- packages beyond the canonical block ----
\usepackage{mathtools}
\usepackage{microtype}
\usepackage{amsthm}
\usepackage[margin=1in]{geometry}
\usepackage{tikz}
\usepackage{tikz-cd}
\usepackage{graphicx}
\usepackage[colorlinks=true,linkcolor=blue!45!black,citecolor=green!35!black,urlcolor=blue!55!black]{hyperref}
\usepackage{cleveref}
\usepackage{xcolor}
\usetikzlibrary{calc}

% ---- paper-local macros (NOT program objects; see coordination/notation.md) ----
% \qA is the quasi-local observable algebra, denoted A in the notation table.
\newcommand{\qA}{\mathcal{A}}
\newcommand{\Aut}{\operatorname{Aut}}
\newcommand{\Der}{\operatorname{Der}}
\newcommand{\id}{\mathrm{id}}
\newcommand{\Ball}[1]{\mathcal{B}_{F}(#1)}      % closed F-norm ball of radius #1
\newcommand{\norm}[1]{\lVert #1\rVert}
\newcommand{\Fnorm}[1]{\lVert #1\rVert_{F}}
\newcommand{\Fanorm}[1]{\lVert #1\rVert_{F_a}}
\newcommand{\loc}{\mathrm{loc}}
\newcommand{\supp}{\operatorname{supp}}
\newcommand{\dist}{\operatorname{dist}}
\newcommand{\R}{\mathbb{R}}
\newcommand{\C}{\mathbb{C}}
\newcommand{\Z}{\mathbb{Z}}
\newcommand{\N}{\mathbb{N}}
\newcommand{\Set}{\mathbf{Cond}}
\newcommand{\lCond}{\mathbf{Cond}^{\mathrm{light}}}
\newcommand{\eqdef}{\coloneqq}

% ---- theorem environments ----
\newtheorem{theorem}{Theorem}[section]
\newtheorem{proposition}[theorem]{Proposition}
\newtheorem{lemma}[theorem]{Lemma}
\newtheorem{corollary}[theorem]{Corollary}
\theoremstyle{definition}
\newtheorem{definition}[theorem]{Definition}
\newtheorem{conjecture}[theorem]{Conjecture}
\newtheorem{example}[theorem]{Example}
\theoremstyle{remark}
\newtheorem{remark}[theorem]{Remark}

% ---- GrokRxiv DOI sidebar (page 1 only) ----
\definecolor{grokgray}{RGB}{110,110,110}
\AddToHook{shipout/foreground}{%
  \ifnum\value{page}=1
    \begin{tikzpicture}[remember picture, overlay]
      \node[
        rotate=90,
        anchor=south,
        font=\Large\sffamily\bfseries\color{grokgray},
        inner sep=0pt
      ] at ([xshift=38pt, yshift=0.52\paperheight]current page.south west)
      {GrokRxiv:2026.07.lieb-robinson-locality\quad
       [\,math-ph\,]\quad
       14 Jul 2026};
    \end{tikzpicture}
  \fi
}

\title{\bfseries Condensed Locality: Lieb--Robinson Estimates and\\
Quasi-Local Dynamics on the Moduli Stack of Hamiltonians}

\author{Matthew Long \\
\textit{The YonedaAI Collaboration} \\
\textit{YonedaAI Research Collective} \\
Chicago, IL \\
\texttt{matthew@yonedaai.com} $\cdot$ \url{https://yonedaai.com}}

\date{14 July 2026\\[4pt]\small Part I of a six-part series on topological phases of matter in the condensed-mathematics paradigm}

\begin{document}
\maketitle

\begin{abstract}
The condensed-mathematics program for topological phases replaces the topological space of
admissible interactions by a condensed object probed by profinite sets, so that continuous
families, profinite disorder hulls, and inverse limits are handled inside a single category
with exact homological algebra. This opening paper supplies the analytic ground floor. Fixing
a Nachtergaele--Sims--Young $\Ffun$-function on a lattice, we recall that the $G$-symmetric
interactions of finite $\Ffun$-norm form a Banach space $\BF$, and we condense it:
$\cond{\BF}(S)=\Cont(S,\BF)$ on profinite probes. We prove three theorems, each with a
complete proof assembled from present-day results. First, $\BF$ is a real Banach space and
the Heisenberg dynamics is a strongly continuous one-parameter group of $*$-automorphisms of
the quasi-local algebra obeying a Lieb--Robinson bound. Second, $\cond{\BF}$ is a condensed
$\R$-vector space and the passage to dynamics is a morphism of condensed sets
$\cond{\R}\times\cond{\BF}\to\cond{\Aut(\qA)}$, by full faithfulness of condensation on
compactly generated spaces; for separable systems the morphism is one of light condensed
sets. Third, an $S$-continuous family of interactions over a profinite base is exactly a
norm-convergent system of locally constant families over finite quotients, and the
load-bearing point is that a uniform $\Ffun$-norm bound forces a Lieb--Robinson velocity and
constants that are uniform over the base. Quasi-adiabatic continuation and
Bachmann--Michalakis--Nachtergaele--Sims automorphic equivalence are then the morphism-level
tools that generate the equivalence class $\Weq$ inverted by the companion papers. We close
with three numbered conjectures: that $\Ham_{d,G}$ is a condensed higher stack, that
time-dependent $\Ffun$-function interactions generate a condensed group with quasi-adiabatic
continuation as internal path-lifting, and that Lieb--Robinson estimates satisfy descent
along finite quotients of a disorder hull, together with a Haskell simulation of the
transverse-field Ising chain that exhibits the light cone and fits the velocity.
\end{abstract}

\tableofcontents

\section{Introduction}
\label{sec:intro}

\subsection{The locality problem in the condensed program}
\label{sec:intro-problem}

A gapped quantum lattice system is specified by a lattice $L$ of sites in dimension $d$, a
finite-dimensional Hilbert space at each site, and a rule that assigns to each finite cluster
of sites a Hermitian operator supported there, an \emph{interaction}. The set of admissible
$G$-symmetric interactions, together with a locality (decay) condition, is the topological
space $\Int_{d,G}$ of the program this series develops. The organizing proposal is to stop
treating $\Int_{d,G}$ as a bare topological space and instead pass to its condensation,
\[
  X \;\longmapsto\; \cond{X}, \qquad \cond{X}(S)=\Cont(S,X),
\]
with $S$ ranging over profinite sets, and then to build over it the moduli stack
$\Ham_{d,G}$ of Hamiltonians, the uniformly gapped substack $\Gap_{d,G}$, the stabilized
phase $\infty$-groupoid $\Phase_{d,G}$, and finally the invertible condensed phase spectrum
$\IPcond_{d,G}$. In one sentence, the thesis of the program is that
\begin{quote}
  a topological phase is a component of the stabilized condensed stack of gapped systems.
\end{quote}
The advantage of this move is not a new Chern number; it is that condensed abelian groups
form a Grothendieck abelian category with exact derived functors, so that continuous
families, profinite disorder, analytic completions, symmetry, stacking, defects, and
transition loci can be manipulated together, where topological groups behave badly.

None of this is available until one controls the analysis at the bottom. Before $\Ham_{d,G}$
can be a condensed object at all, three things must be true and must be \emph{uniform in the
probe}: the space of interactions must be a Banach space so that its condensation is a
condensed $\R$-vector space; the Heisenberg dynamics must depend continuously on the
interaction and the time in the Banach topology, so that a family of dynamics is a morphism
rather than a set-theoretic assignment; and the Lieb--Robinson estimates that make the
dynamics quasi-local must survive the passage to profinite limits with constants that do not
degenerate. This paper establishes exactly these facts, and no more. The word ``uniform'' is
doing real work throughout: a family of quasi-local dynamics whose Lieb--Robinson velocity is
finite pointwise but unbounded over the base is not a morphism of condensed sets, just as a
family whose gap is positive pointwise but has infimum zero is not a member of $\Gap_{d,G}$.

The tools are classical. Lieb and Robinson proved a finite group velocity for quantum spin
systems in 1972~\cite{lieb-robinson-1972}; the modern reformulation through
$\Ffun$-functions, due to Nachtergaele, Sims, and Young~\cite{nsy-quasilocality-1,
nsy-quasilocality-2} and surveyed by Hastings~\cite{hastings-locality}, gives exactly the
Banach-space packaging we need. The condensed formalism is that of Clausen and
Scholze~\cite{clausen-scholze-condensed, clausen-scholze-analytic,
clausen-scholze-analytic-stacks}, with the pyknotic variant of Barwick and
Haine~\cite{barwick-haine-pyknotic}. The closest precedent on the physics side is the
homotopical, sheaf-theoretic study of parametrized families of spin systems by Beaudry and
collaborators~\cite{beaudry-etal}, which organizes families of gapped systems with
ordinary, not condensed, topology; the present paper is the condensed refinement of that
viewpoint at the locality layer. Our contribution is to show that these two mature bodies
of work fit together on the nose: the analytic estimate that Nachtergaele--Sims--Young prove
is precisely the continuity statement that condensation requires, and the profinite probes of
condensed mathematics are precisely matched to the inverse limits of finite-volume and
disorder data that lattice physics produces.

\subsection{What this paper proves, and what it only conjectures}
\label{sec:intro-results}

We adopt the editorial stance of the whole series: a statement is labelled
Theorem or Proposition only when it has a complete proof from cited present-day results, and
everything past that boundary is a numbered Conjecture. Under that discipline the results are
as follows. \Cref{thm:banach} (labelled I-A) records that the $\Ffun$-normed interactions
form a real Banach space $\BF$ and that the Heisenberg dynamics is a strongly continuous
one-parameter group of $*$-automorphisms with a Lieb--Robinson bound; this is a repackaging
of~\cite{nsy-quasilocality-1}. \Cref{thm:condensed-dynamics} (I-B) upgrades this to condensed
mathematics: $\cond{\BF}$ is a condensed $\R$-vector space, and the assignment
$(t,\Phi)\mapsto\tau^\Phi_t$ is a morphism of condensed sets, because it is continuous in the
Banach topology and condensation is fully faithful on compactly generated spaces. \Cref{thm:uniform-lr}
(I-C) is the analytic heart: over a profinite base, an $S$-continuous family of interactions
is the same datum as a norm-convergent compatible system over the finite quotients, and a
uniform bound on the $\Ffun$-norm produces a Lieb--Robinson light cone whose velocity and
prefactors are uniform in $s\in S$. \Cref{prop:clustering} draws the uniform exponential
clustering consequence that Part III will need, and \Cref{prop:qac} records that
quasi-adiabatic continuation and automorphic equivalence~\cite{hastings-wen-qac,
bmns-automorphic} furnish the morphism-level generators of the equivalence class $\Weq$.

The conjectures mark the edge of what today's estimates deliver. \Cref{conj:stack} (I-1), that
$\Ham_{d,G}$ is a condensed \emph{higher} stack (the groupoid of gauge and quasi-local
automorphisms glues along finite covers of profinite probes), is not proved here and is not
asserted anywhere in the series as a theorem. \Cref{conj:group} (I-2) proposes that the
quasi-local automorphisms generated by time-dependent $\Ffun$-function interactions form a
condensed group for which quasi-adiabatic continuation is an internal path-lifting.
\Cref{conj:descent} (I-3) proposes that Lieb--Robinson estimates satisfy descent along finite
quotients of a disorder hull $\dis=Q^{\Z^d}$, so that quasi-local invariants are
determined by finite-resolution data. These three are the locality-layer instances of the
program's Master Conjecture.

\subsection{Relation to companion papers}
\label{sec:intro-companions}

This is Part~I of six. The series is modular: each part takes the previous ones as input and
produces new structure, rather than assembling into a single monolith.

Part~II, \emph{Positivity, $C^*$-Norms, and Condensed State Spaces of Quasi-Local
Algebras}~\cite{paperPositivity}, takes the quasi-local algebra $\qA$ and the automorphism
morphism built here and attaches the condensed state space and the solid-$K$-theory invariant
of Aoki~\cite{aoki-solidification}. The Lieb--Robinson bound of \Cref{thm:banach} is what
makes the Heisenberg dynamics act on $\qA$ by genuine $*$-automorphisms, which Part~II needs
before it can speak of states and positivity. Part~III, \emph{The Uniformly Gapped
Substack}~\cite{paperGap}, cuts out $\Gap_{d,G}\subset\Ham_{d,G}$ and studies stability of the gap;
its exponential-clustering input is our \Cref{prop:clustering}, and its automorphic
equivalences are the $\Weq$ generated by our \Cref{prop:qac}. Part~IV, \emph{From Lattice
Models to Effective Field Theories}~\cite{paperEFT}, group-completes the invertible sector
into $\IPcond_{d,G}$; the continuity of dynamics in the interaction (\Cref{thm:condensed-dynamics})
is what lets its parametrized families be morphisms of condensed objects. Part~V,
\emph{Physical Realizability of Bordism and Homotopy Classes}~\cite{paperRealizability},
asks which abstract classes are realized by uniformly gapped lattice families, and inherits
the uniform-over-the-base discipline established here. Part~VI~\cite{paperSynthesis} composes
the five modules and states the Master Conjecture. Where those papers speak of ``an
$S$-continuous family of Hamiltonians'', the precise meaning is the one fixed by
\Cref{thm:uniform-lr}.

\medskip
The rest of the paper is organized as follows. \Cref{sec:framework} fixes the lattice, the
quasi-local algebra, the $\Ffun$-functions, and the Banach space $\BF$. \Cref{sec:lr} recalls
the Heisenberg dynamics and the Lieb--Robinson bound and proves I-A. \Cref{sec:condensation}
recalls exactly as much condensed mathematics as we use and proves the condensed vector-space
and dynamics statements (I-B). \Cref{sec:uniform} proves the uniform Lieb--Robinson bound over
a profinite base and the compatible-finite-data description (I-C), and draws the clustering
corollary. \Cref{sec:qac} treats quasi-adiabatic continuation and the generation of $\Weq$.
\Cref{sec:conjectures} states the three conjectures. \Cref{sec:numerics} reports the
transverse-field Ising simulation. \Cref{sec:discussion,sec:conclusion} discuss limitations
and conclude.

\section{Quantum lattice systems and interaction spaces}
\label{sec:framework}

\subsection{Lattice, local Hilbert spaces, and the quasi-local algebra}
\label{sec:framework-lattice}

We work in the standard operator-algebraic setting for quantum spin systems; the reader who
wants proofs of the folklore statements below will find them in~\cite{nsy-quasilocality-1}
and the references there.

\begin{definition}[Lattice and local structure]
\label{def:lattice}
A \emph{lattice} is a countable set $L$ equipped with a metric $\dist$ such that the balls
$b_x(r)=\{y\in L : \dist(x,y)\le r\}$ are finite and their cardinality grows at most
polynomially: there are $\kappa,\nu>0$ with $\lvert b_x(r)\rvert\le\kappa(1+r)^\nu$ for all
$x\in L$ and $r\ge 0$. The integer lattice $\Z^d$ with the $\ell^1$ metric is the running
example, with $\nu=d$. To each site $x\in L$ we attach a finite-dimensional Hilbert space
$\mathcal H_x$ of dimension $n_x$, uniformly bounded: $\sup_x n_x<\infty$.
\end{definition}

For a finite $X\subset L$ write $\mathcal H_X=\bigotimes_{x\in X}\mathcal H_x$ and let
$\qA_X=\mathcal B(\mathcal H_X)$ be the matrix algebra of observables supported on $X$. If
$X\subset Y$ then $\qA_X\hookrightarrow\qA_Y$ by $A\mapsto A\otimes\mathbf 1_{Y\setminus X}$,
and the \emph{local algebra} is the union
\[
  \qA_\loc \;=\; \bigcup_{X\Subset L}\qA_X,
\]
the colimit over finite subsets $X\Subset L$. Its completion in the operator norm is the
\emph{quasi-local algebra}
\[
  \qA \;=\; \overline{\qA_\loc}^{\,\norm{\cdot}},
\]
a unital $C^*$-algebra; this is the standard $C^*$-inductive-limit (quasi-local) construction
of Bratteli--Robinson~\cite{bratteli-robinson-1}. Because $L$ is countable and each $\qA_X$ is
finite-dimensional, $\qA$ is separable. An element $A\in\qA$ is \emph{local} if $A\in\qA_X$ for some finite $X$,
and its \emph{support} $\supp A$ is the smallest such $X$. We write $\qA$ for what the
notation table of the series calls the observable algebra $A$; the calligraphic letter is the
standard one in the Lieb--Robinson literature and we keep it to avoid clashing with generic
operators.

A symmetry is a group $G$ acting on $\qA$ by $*$-automorphisms, on-site in the internal case
(a representation $u_x:G\to\mathcal U(\mathcal H_x)$ inducing $\alpha_g=\bigotimes_x
\mathrm{Ad}(u_x(g))$) and by isometries of $(L,\dist)$ in the spatial case. Everything below
is compatible with a fixed such $G$; to keep the analysis in the foreground we suppress $G$
from the notation and reinstate it only where symmetry changes a statement.

\subsection{\texorpdfstring{$\Ffun$}{F}-functions and the interaction Banach space}
\label{sec:framework-Ffun}

The decay of interactions is measured against a weight on the metric, an $\Ffun$-function.

\begin{definition}[$\Ffun$-function]
\label{def:Ffun}
An \emph{$\Ffun$-function} on $(L,\dist)$ is a non-increasing function
$\Ffun:[0,\infty)\to(0,\infty)$ satisfying
\begin{enumerate}
\item[(F1)] \emph{Uniform summability:}
  $\displaystyle \norm{\Ffun}\eqdef\sup_{x\in L}\sum_{y\in L}\Ffun\bigl(\dist(x,y)\bigr)<\infty$;
\item[(F2)] \emph{Convolution bound:}
  $\displaystyle C_\Ffun\eqdef\sup_{x,y\in L}\;\sum_{z\in L}
   \frac{\Ffun(\dist(x,z))\,\Ffun(\dist(z,y))}{\Ffun(\dist(x,y))}<\infty.$
\end{enumerate}
\end{definition}

On $\Z^d$ the functions $\Ffun_\varepsilon(r)=(1+r)^{-(d+\varepsilon)}$ with $\varepsilon>0$
are $\Ffun$-functions, as are the exponentially weighted $\Ffun(r)=e^{-\theta r}
(1+r)^{-(d+\varepsilon)}$. The next lemma, which we use constantly, says that one may always
sharpen the decay of a given $\Ffun$-function by an exponential factor at no cost in the two
constants; it is the mechanism by which a finite velocity appears.

\begin{lemma}[Exponential reweighting]
\label{lem:reweight}
Let $\Ffun$ be an $\Ffun$-function and $a\ge 0$. Then $\Ffun_a(r)\eqdef e^{-ar}\Ffun(r)$ is
an $\Ffun$-function with
\[
  \norm{\Ffun_a}\le\norm{\Ffun}, \qquad C_{\Ffun_a}\le C_\Ffun .
\]
\end{lemma}
\begin{proof}
$\Ffun_a$ is positive and non-increasing. For (F1), $e^{-ar}\le 1$ gives
$\Ffun_a(r)\le\Ffun(r)$, so $\norm{\Ffun_a}\le\norm{\Ffun}$. For (F2), the triangle
inequality $\dist(x,y)\le\dist(x,z)+\dist(z,y)$ and monotonicity of $t\mapsto e^{-at}$ give
$e^{-a\dist(x,z)}e^{-a\dist(z,y)}=e^{-a(\dist(x,z)+\dist(z,y))}\le e^{-a\dist(x,y)}$, hence
\begin{multline*}
  \frac{\Ffun_a(\dist(x,z))\,\Ffun_a(\dist(z,y))}{\Ffun_a(\dist(x,y))}
  =e^{-a(\dist(x,z)+\dist(z,y)-\dist(x,y))}\,
   \frac{\Ffun(\dist(x,z))\,\Ffun(\dist(z,y))}{\Ffun(\dist(x,y))}\\
  \le \frac{\Ffun(\dist(x,z))\,\Ffun(\dist(z,y))}{\Ffun(\dist(x,y))}.
\end{multline*}
Summing over $z$ and taking the supremum over $x,y$ yields $C_{\Ffun_a}\le C_\Ffun$.
\end{proof}

To keep the constants concrete---the knowledge base for this program insists that the
$\Ffun$-norm never be left implicit---we verify the axioms for the standard polynomial weights.

\begin{lemma}[Polynomial $\Ffun$-functions on $\Z^d$]
\label{lem:zd-Ffun}
On $(\Z^d,\ell^1)$ the function $\Ffun_\varepsilon(r)=(1+r)^{-(d+\varepsilon)}$ is an
$\Ffun$-function for every $\varepsilon>0$, with $\norm{\Ffun_\varepsilon}<\infty$ and
$C_{\Ffun_\varepsilon}\le 2^{\,d+\varepsilon+1}\,\norm{\Ffun_\varepsilon}$, both independent of
the base point by translation invariance.
\end{lemma}
\begin{proof}
Throughout, $\lvert x-y\rvert_1=\dist(x,y)$ denotes the $\ell^1$ distance and
$\lvert y\rvert_1=\dist(0,y)$ the distance to the origin.
The number of sites at $\ell^1$-distance $n$ from a fixed point is $O(n^{d-1})$, so
\[
  \norm{\Ffun_\varepsilon}=\sum_{y\in\Z^d}(1+\lvert y\rvert_1)^{-(d+\varepsilon)}
  \le c_d\sum_{n\ge 0}(1+n)^{d-1}(1+n)^{-(d+\varepsilon)}
  = c_d\sum_{n\ge 0}(1+n)^{-(1+\varepsilon)}<\infty .
\]
For (F2), split the sum over $z$ according to whether $\lvert x-z\rvert_1\ge\tfrac12\lvert
x-y\rvert_1$ or $\lvert z-y\rvert_1\ge\tfrac12\lvert x-y\rvert_1$; at least one holds by the
triangle inequality. In the first case $(1+\lvert x-y\rvert_1)/(1+\lvert x-z\rvert_1)\le 2$, so
$\Ffun_\varepsilon(\lvert x-z\rvert_1)/\Ffun_\varepsilon(\lvert x-y\rvert_1)\le
2^{\,d+\varepsilon}$ and the summand $\Ffun_\varepsilon(\lvert x-z\rvert_1)\Ffun_\varepsilon
(\lvert z-y\rvert_1)/\Ffun_\varepsilon(\lvert x-y\rvert_1)$ is at most
$2^{\,d+\varepsilon}\Ffun_\varepsilon(\lvert z-y\rvert_1)$, whose sum over $z$ is
$2^{\,d+\varepsilon}\norm{\Ffun_\varepsilon}$. The second case is symmetric, giving
$C_{\Ffun_\varepsilon}\le 2^{\,d+\varepsilon+1}\norm{\Ffun_\varepsilon}<\infty$.
\end{proof}

\begin{definition}[Interaction and $\Ffun$-norm]
\label{def:interaction}
An \emph{interaction} is a map $\Phi$ that assigns to each finite $X\Subset L$ a
self-adjoint element $\Phi(X)=\Phi(X)^*\in\qA_X$, and in the $G$-symmetric case commutes with
the symmetry, $\alpha_g(\Phi(X))=\Phi(gX)$ for all $g\in G$. Its \emph{$\Ffun$-norm} is
\[
  \Fnorm{\Phi}\;\eqdef\;\sup_{x,y\in L}\;\frac{1}{\Ffun(\dist(x,y))}
    \sum_{X\ni x,y}\norm{\Phi(X)},
\]
the inner sum ranging over finite $X$ containing both $x$ and $y$. The
\emph{interaction Banach space} is
\[
  \BF\;\eqdef\;\bigl\{\,\Phi \text{ an interaction} : \Fnorm{\Phi}<\infty\,\bigr\}.
\]
\end{definition}

The $\Ffun$-norm controls, uniformly in the site, the total weight of interaction terms that
link any two sites. A finite-range interaction of bounded strength has finite $\Ffun$-norm
for every $\Ffun$-function; a two-body interaction $\Phi(\{x,y\})$ decaying like
$\Ffun(\dist(x,y))$ does too. The point of \Cref{def:interaction} is that everything downstream
depends on $\Phi$ only through the single number $\Fnorm{\Phi}$ (and the constants of the
chosen $\Ffun$), which is what will let us pass to families.

\begin{proposition}[$\BF$ is a Banach space]
\label{prop:banach-space}
$\bigl(\BF,\Fnorm{\cdot}\bigr)$ is a real Banach space, and the map
$\Phi\mapsto\Fnorm{\Phi}$ is a genuine norm on it.
\end{proposition}
\begin{proof}
Absolute homogeneity and the triangle inequality are inherited from the operator norm inside
the supremum and the sum, and $\Fnorm{\Phi}=0$ forces $\Phi(X)=0$ for every $X$ (take
$x,y\in X$). Only completeness needs an argument. Let $(\Phi^{(n)})_n$ be Cauchy in
$\Fnorm{\cdot}$. For each fixed finite $X$ and any two sites $x,y\in X$,
\[
  \norm{\Phi^{(n)}(X)-\Phi^{(m)}(X)}
  \le \Ffun(\dist(x,y))\,\Fnorm{\Phi^{(n)}-\Phi^{(m)}},
\]
so $(\Phi^{(n)}(X))_n$ is Cauchy in the finite-dimensional space $\qA_X$ and converges to some
self-adjoint $\Phi(X)\in\qA_X$; the limit is symmetric if each $\Phi^{(n)}$ is. Fix
$\varepsilon>0$ and $N$ with $\Fnorm{\Phi^{(n)}-\Phi^{(m)}}\le\varepsilon$ for $n,m\ge N$. For
every pair $x,y$ and every finite truncation of the sum $\sum_{X\ni x,y}$, letting
$m\to\infty$ gives
$\Ffun(\dist(x,y))^{-1}\sum_{X\ni x,y}\norm{\Phi^{(n)}(X)-\Phi(X)}\le\varepsilon$ for $n\ge N$,
uniformly in $x,y$; hence $\Fnorm{\Phi^{(n)}-\Phi}\le\varepsilon$ and $\Phi\in\BF$ with
$\Phi^{(n)}\to\Phi$. Thus $\BF$ is complete.
\end{proof}

We record the two elementary structures that condensation will act on. The self-adjointness
constraint makes $\BF$ a real (not complex) Banach space; scaling an interaction by a real
number and adding interactions termwise are continuous, so $\BF$ is in particular a
topological real vector space and an abelian topological group under addition.

\section{The Heisenberg dynamics and the Lieb--Robinson bound}
\label{sec:lr}

\subsection{Finite-volume dynamics and the infinite-volume limit}
\label{sec:lr-dynamics}

For a finite $\Lambda\Subset L$ the \emph{local Hamiltonian} of $\Phi\in\BF$ is the
self-adjoint operator
\[
  H^\Phi_\Lambda \;=\; \sum_{X\subseteq\Lambda}\Phi(X)\;\in\;\qA_\Lambda,
\]
a finite sum, and the finite-volume \emph{Heisenberg dynamics} is the one-parameter group of
$*$-automorphisms
\[
  \tau^{\Phi,\Lambda}_t(A)\;=\;e^{itH^\Phi_\Lambda}\,A\,e^{-itH^\Phi_\Lambda},\qquad A\in\qA.
\]
The Lieb--Robinson bound controls the spatial spread of $\tau^{\Phi,\Lambda}_t(A)$ and is the
key to removing the cutoff $\Lambda$.

\begin{theorem}[Lieb--Robinson bound, $\Ffun$-function form~\cite{nsy-quasilocality-1}]
\label{thm:lr-bound}
Let $\Phi\in\BF$ and let $X,Y\Subset L$ be disjoint. For $A\in\qA_X$, $B\in\qA_Y$, and every
finite $\Lambda\supseteq X\cup Y$,
\[
  \bigl\lVert[\tau^{\Phi,\Lambda}_t(A),B]\bigr\rVert
  \;\le\;\frac{2\norm{A}\,\norm{B}}{C_\Ffun}\,
    \Bigl(e^{2\Fnorm{\Phi}\,C_\Ffun\,\lvert t\rvert}-1\Bigr)
    \sum_{x\in X}\sum_{y\in Y}\Ffun\bigl(\dist(x,y)\bigr),
\]
with $\norm{\Ffun}$ and $C_\Ffun$ as in \Cref{def:Ffun}. The bound is uniform in $\Lambda$.
\end{theorem}

This is Theorem~3.4 of~\cite{nsy-quasilocality-1} in the present notation; we take it as
given. The velocity is made explicit by combining it with reweighting.

\begin{proposition}[Light cone and velocity]
\label{prop:velocity}
Let $\Phi\in\mathcal B_{\Ffun_a}$ for the reweighted function $\Ffun_a(r)=e^{-ar}\Ffun(r)$ of
\Cref{lem:reweight}, with $a>0$. Then for disjoint $X,Y$ and $A\in\qA_X$, $B\in\qA_Y$,
\[
  \bigl\lVert[\tau^{\Phi,\Lambda}_t(A),B]\bigr\rVert
  \;\le\;\frac{2\norm{A}\,\norm{B}\,\norm{\Ffun}}{C_{\Ffun_a}}\,
    \min\{\lvert X\rvert,\lvert Y\rvert\}\;
    e^{-a\bigl(\dist(X,Y)-v\lvert t\rvert\bigr)},
    \qquad v \;=\; \frac{2\,\Fanorm{\Phi}\,C_{\Ffun_a}}{a},
\]
where $\dist(X,Y)=\min_{x\in X,y\in Y}\dist(x,y)$. In particular the commutator is
exponentially small outside the light cone $\dist(X,Y)\ge v\lvert t\rvert$.
\end{proposition}
\begin{proof}
Apply \Cref{thm:lr-bound} with $\Ffun_a$ in place of $\Ffun$. For the geometric sum, each
$x\in X$ has $\sum_{y\in Y}\Ffun_a(\dist(x,y))\le e^{-a\dist(X,Y)}\sum_{y\in
L}\Ffun(\dist(x,y))\le\norm{\Ffun}e^{-a\dist(X,Y)}$ using $e^{-a\dist(x,y)}\le
e^{-a\dist(X,Y)}$; summing over the smaller of $X,Y$ gives the factor
$\min\{\lvert X\rvert,\lvert Y\rvert\}$. Finally $e^{2\Fanorm{\Phi}C_{\Ffun_a}\lvert
t\rvert}-1\le e^{2\Fanorm{\Phi}C_{\Ffun_a}\lvert t\rvert}=e^{av\lvert t\rvert}$, and the two
exponentials combine.
\end{proof}

\begin{theorem}[Banach interaction space and strongly continuous LR dynamics, I-A]
\label{thm:banach}
Fix an $\Ffun$-function $\Ffun$ and $\Phi\in\BF$.
\begin{enumerate}
\item $\BF$ is a real Banach space (\Cref{prop:banach-space}).
\item For each $A\in\qA$ the limit
  $\tau^\Phi_t(A)=\lim_{\Lambda\uparrow L}\tau^{\Phi,\Lambda}_t(A)$
  exists in operator norm, uniformly for $t$ in compact sets, and defines a strongly
  continuous one-parameter group $\{\tau^\Phi_t\}_{t\in\R}$ of $*$-automorphisms of $\qA$.
\item The infinite-volume dynamics satisfies the Lieb--Robinson bound of
  \Cref{thm:lr-bound} and the light-cone estimate of \Cref{prop:velocity} with the same
  constants (the finite-volume bounds being uniform in $\Lambda$).
\end{enumerate}
\end{theorem}
\begin{proof}
Part~(1) is \Cref{prop:banach-space}. For part~(2): for local $A$ and finite volumes
$\Lambda\subseteq\Lambda'$, the difference $\tau^{\Phi,\Lambda'}_t(A)-\tau^{\Phi,\Lambda}_t(A)$
is expressed by Duhamel's formula as a time integral of a commutator of
$H^\Phi_{\Lambda'}-H^\Phi_\Lambda$ with $\tau^{\Phi,\Lambda}_{s}(A)$; the Lieb--Robinson bound
of \Cref{thm:lr-bound} bounds that commutator by the tail $\sum_{X\ni x,\,X\not\subseteq
\Lambda}\norm{\Phi(X)}$ against $\Ffun$, which is a convergent tail of a $\Fnorm{\cdot}$-finite
sum and hence tends to $0$ as $\Lambda\uparrow L$, uniformly for $t$ in compacts. This is the
Cauchy criterion, giving a limit $\tau^\Phi_t(A)$ for local $A$; the estimate is uniform in
$A$ on norm-balls, so the limit extends to all of $\qA=\overline{\qA_\loc}$ by density and is
an isometric $*$-homomorphism. The group law and strong continuity pass to the limit. This is
the content of~\cite[\S3--4]{nsy-quasilocality-1}; see also~\cite{hastings-locality} for the
argument in the exponential case and~\cite{lieb-robinson-1972} for the original velocity.
Part~(3) is the uniformity in $\Lambda$ of \Cref{thm:lr-bound,prop:velocity}, which survives
the limit.
\end{proof}

The proof of \Cref{thm:banach}(2) already contains the estimate we need for the condensed
statement, because it is really a statement about how $\tau$ depends on $\Phi$. We isolate it.

\begin{lemma}[Continuity of the dynamics in the interaction]
\label{lem:dynamics-lipschitz}
Fix a local observable $A\in\qA_X$ and a time horizon $T>0$. On the closed ball
$\Ball{B}=\{\Phi:\Fnorm{\Phi}\le B\}$ there is a constant
$K=K(A,T,B,\Ffun)<\infty$ such that
\[
  \sup_{\lvert t\rvert\le T}\bigl\lVert\tau^\Phi_t(A)-\tau^\Psi_t(A)\bigr\rVert
  \;\le\;K\,\Fnorm{\Phi-\Psi}
  \qquad\text{for all }\Phi,\Psi\in\Ball{B}.
\]
Moreover $t\mapsto\tau^\Phi_t(A)$ is norm-continuous, so
$(t,\Phi)\mapsto\tau^\Phi_t(A)$ is jointly continuous $\R\times\Ball{B}\to\qA$.
\end{lemma}
\begin{proof}
Work first in a finite volume $\Lambda\supseteq X$ and interpolate along the straight line
$\Phi_r=\Psi+r(\Phi-\Psi)$, $r\in[0,1]$. Duhamel gives
\[
  \tau^{\Phi,\Lambda}_t(A)-\tau^{\Psi,\Lambda}_t(A)
  = i\int_0^t \tau^{\Phi,\Lambda}_{s}
      \bigl(\,[\,H^{\Phi}_\Lambda-H^{\Psi}_\Lambda,\;
              \tau^{\Psi,\Lambda}_{t-s}(A)\,]\,\bigr)\,ds .
\]
The generator difference is $H^{\Phi}_\Lambda-H^{\Psi}_\Lambda=\sum_{Z\subseteq\Lambda}
(\Phi-\Psi)(Z)$. Bounding the commutator of each term $(\Phi-\Psi)(Z)$ with the
Lieb--Robinson-evolved $\tau^{\Psi,\Lambda}_{t-s}(A)$ by \Cref{thm:lr-bound}, and summing the
resulting weights against $\Ffun$, produces
\[
  \bigl\lVert[\,H^{\Phi}_\Lambda-H^{\Psi}_\Lambda,\;\tau^{\Psi,\Lambda}_{t-s}(A)\,]\bigr\rVert
  \;\le\; 2\norm{A}\,\Fnorm{\Phi-\Psi}\,\lvert X\rvert\,\norm{\Ffun}\,
     e^{2BC_\Ffun\lvert t-s\rvert},
\]
where the exponential comes from the Lieb--Robinson factor at $\Fnorm{\Psi}\le B$ and the
factor $\lvert X\rvert\norm{\Ffun}$ from summing $\Ffun$ over the support of $A$ against all
sites. Integrating over $\lvert s\rvert\le\lvert t\rvert\le T$ gives the claim in finite
volume with $K=2\norm{A}\lvert X\rvert\norm{\Ffun}\,T\,e^{2BC_\Ffun T}$, uniform in $\Lambda$;
letting $\Lambda\uparrow L$ using \Cref{thm:banach}(2) gives it for the infinite-volume
dynamics. Norm-continuity in $t$ is the strong continuity of \Cref{thm:banach}(2). Joint
continuity follows from the two one-variable estimates and the triangle inequality.
\end{proof}

\subsection{The generator as a symmetric derivation}
\label{sec:lr-generator}

The strongly continuous group $\tau^\Phi$ has an infinitesimal generator, and the
$\Ffun$-norm controls it on local observables. This is the infinitesimal counterpart of the
whole picture: the same single number $\Fnorm{\Phi}$ that bounds the dynamics also bounds its
generator, term by term.

\begin{proposition}[Symmetric derivation generating the dynamics]
\label{prop:generator}
Fix $\Phi\in\BF$. For local $A\in\qA_X$ the sum
\[
  \delta_\Phi(A)\;=\;i\sum_{Z:\,Z\cap X\ne\emptyset}[\,\Phi(Z),A\,]
\]
converges in operator norm, with
$\norm{\delta_\Phi(A)}\le 2\,\lvert X\rvert\,\Ffun(0)\,\Fnorm{\Phi}\,\norm{A}$, and defines a
symmetric derivation $\delta_\Phi$ on the dense $*$-subalgebra $\qA_\loc$. Its closure
generates $\tau^\Phi$, i.e.\ $\tau^\Phi_t=\exp(t\,\overline{\delta_\Phi})$ as strongly
continuous groups, and $\tfrac{d}{dt}\big|_{t=0}\tau^\Phi_t(A)=\delta_\Phi(A)$ for local $A$.
\end{proposition}
\begin{proof}
For $A\in\qA_X$ only clusters $Z$ meeting $X$ contribute. Bounding each commutator by
$2\norm{\Phi(Z)}\norm{A}$ and using $\sum_{Z\ni x}\norm{\Phi(Z)}\le\Ffun(0)\Fnorm{\Phi}$ for
each $x\in X$---this is the $\Ffun$-norm evaluated at the diagonal pair $y=x$, since
$\Ffun(0)^{-1}\sum_{Z\ni x}\norm{\Phi(Z)}\le\Fnorm{\Phi}$---gives
\[
  \norm{\delta_\Phi(A)}\;\le\;2\norm{A}\sum_{x\in X}\sum_{Z\ni x}\norm{\Phi(Z)}
  \;\le\;2\,\lvert X\rvert\,\Ffun(0)\,\Fnorm{\Phi}\,\norm{A} ,
\]
so the sum converges absolutely. The Leibniz rule $\delta_\Phi(AB)=\delta_\Phi(A)B+A\delta_\Phi
(B)$ and the symmetry $\delta_\Phi(A^*)=\delta_\Phi(A)^*$ are inherited from the commutator, so
$\delta_\Phi$ is a symmetric derivation on $\qA_\loc$. That its closure generates $\tau^\Phi$
and that the group differentiates to $\delta_\Phi$ on local elements is the standard generation
theorem for quasi-local dynamics~\cite{nsy-quasilocality-1, hastings-locality}; the
Lieb--Robinson bound of \Cref{thm:banach} is exactly what makes the finite-volume generators
$\delta^\Lambda_\Phi(A)=i[H^\Phi_\Lambda,A]$ converge to $\delta_\Phi(A)$ as $\Lambda\uparrow L$.
\end{proof}

The assignment $\Phi\mapsto\delta_\Phi$ is real-linear and, on each $\qA_X$, bounded in the
$\Ffun$-norm: $\norm{\delta_\Phi(A)-\delta_\Psi(A)}\le 2\lvert X\rvert\Ffun(0)\Fnorm{\Phi-\Psi}
\norm{A}$. This is the infinitesimal shadow of \Cref{lem:dynamics-lipschitz}; once condensed it
says that the passage to generators is itself a morphism $\cond{\BF}\to\cond{\Der(\qA_\loc)}$,
continuous in the $\Ffun$-norm, of which \Cref{thm:condensed-dynamics} is the exponentiated
form.

\begin{example}[Transverse-field Ising chain]
\label{ex:tfim}
Take $L=\Z$, $\mathcal H_x=\C^2$, and
\[
  \Phi_{\mathrm{TFIM}}(\{x,x+1\})=-J\,\sigma^z_x\sigma^z_{x+1},\qquad
  \Phi_{\mathrm{TFIM}}(\{x\})=-h\,\sigma^x_x,
\]
all other $\Phi(X)=0$. This is finite range, so with any $\Ffun$-function---for instance
$\Ffun_\varepsilon$ of \Cref{lem:zd-Ffun}---the $\Ffun$-norm is finite. The heaviest diagonal
pair $x=y$ collects the on-site field together with the two incident bonds, contributing
$(\lvert h\rvert+2\lvert J\rvert)/\Ffun(0)$; a nearest-neighbour pair collects the single bond
between them, contributing $\lvert J\rvert/\Ffun(1)$; and no pair at distance $\ge 2$
contributes. Since $\Ffun$ is non-increasing, $\Ffun(1)\le\Ffun(0)$, so
\[
  \Fnorm{\Phi_{\mathrm{TFIM}}}
  =\max\!\Bigl\{\tfrac{\lvert h\rvert+2\lvert J\rvert}{\Ffun(0)},\;
     \tfrac{\lvert J\rvert}{\Ffun(1)}\Bigr\}
  \;\le\;\frac{2\lvert J\rvert+\lvert h\rvert}{\Ffun(1)}\;<\;\infty,
\]
and \Cref{thm:banach} applies for all $J,h$. The velocity bound of \Cref{prop:velocity} is
$v=2\Fanorm{\Phi}C_{\Ffun_a}/a$; its optimization over the reweighting parameter $a$ produces
a finite group velocity, which we exhibit numerically in \Cref{sec:numerics}. The chain is
gapped except on the critical line $\lvert J\rvert=\lvert h\rvert$; in the language of the
program the critical line is a slice of the gapless discriminant $\Sig_f$, and the two gapped
phases sit in different components of $\Phase_{d,G}$. The present paper says nothing about the
gap---that is Part~III---but it does guarantee that the dynamics on both sides has the same
uniform light cone as long as $J,h$ stay bounded.
\end{example}

\section{Condensation of the interaction space}
\label{sec:condensation}

\subsection{Condensed sets and profinite probes}
\label{sec:condensation-recall}

We recall only what we use. A \emph{profinite set} is a cofiltered limit $S=\varprojlim_i
S_i$ of finite sets, equivalently a compact, Hausdorff, totally disconnected topological
space. The \emph{site of profinite sets} has covers the finite jointly surjective families;
a \emph{condensed set} is a sheaf on this site (equivalently, on the subcategory of
extremally disconnected profinite sets, where sheaf and product-preserving presheaf coincide).
Condensed abelian groups form a Grothendieck abelian category, so kernels, cokernels, and
derived functors behave as in ordinary homological algebra~\cite{clausen-scholze-condensed};
the pyknotic formulation is~\cite{barwick-haine-pyknotic}. Write $\Set$ for the category of
condensed sets.

\begin{definition}[Condensation]
\label{def:condensation}
For a topological space $X$, its \emph{condensation} $\cond{X}\in\Set$ is
\[
  \cond{X}(S)\;=\;\Cont(S,X),\qquad S\ \text{profinite},
\]
the set of continuous maps $S\to X$, with the evident restriction along maps of profinite
sets. For a finite set $S$ this is just $X^{S}$, so $\cond{X}(\ast)=X$ recovers the points.
\end{definition}

The one structural fact we lean on is full faithfulness. A topological space is
\emph{compactly generated} if its topology is final for the maps out of compact Hausdorff
spaces into it; metrizable spaces and, more generally, first-countable spaces are compactly
generated.

\begin{theorem}[Full faithfulness of condensation~\cite{clausen-scholze-condensed}]
\label{thm:full-faithful}
The functor $X\mapsto\cond{X}$ from compactly generated Hausdorff spaces to condensed sets is
fully faithful and preserves finite products. In particular, for compactly generated
Hausdorff $X,Y$,
\[
  \Cont(X,Y)\;\xrightarrow{\ \sim\ }\;\operatorname{Hom}_{\Set}(\cond{X},\cond{Y}),
\]
so a continuous map is the same datum as a morphism of condensed sets, and
$\cond{X\times Y}\cong\cond{X}\times\cond{Y}$.
\end{theorem}

Two comments fix the size of the formalism we need. First, a Banach space $V$ is metrizable,
hence compactly generated, so \Cref{thm:full-faithful} applies to it; moreover $\cond{V}$ is
a module over the condensed ring $\cond{\R}$, i.e.\ a \emph{condensed $\R$-vector space},
because addition $V\times V\to V$ and scaling $\R\times V\to V$ are continuous and condensation
preserves finite products. Second, for separable target spaces one may replace profinite
probes by their \emph{light} (metrizable, equivalently second-countable) counterparts and work
in the category $\lCond$ of light condensed sets~\cite{clausen-scholze-analytic-stacks}; this
avoids the set-theoretic subtleties of the full theory (no inaccessible cardinal is needed)
and is the correct home for separable spin systems, whose observable algebra $\qA$ is
separable by \Cref{def:lattice}.

\subsection{The condensed interaction space and the condensed dynamics}
\label{sec:condensation-main}

\begin{proposition}[Condensed structure of $\BF$]
\label{prop:condensed-BF}
$\cond{\BF}$ is a condensed real vector space and a condensed abelian group. For a profinite
set $S$,
\[
  \cond{\BF}(S)=\Cont(S,\BF)
\]
is the real Banach space of continuous $\BF$-valued functions on $S$ with the supremum norm
$\norm{\Phi_\bullet}_{\infty,\Ffun}=\sup_{s\in S}\Fnorm{\Phi_s}$, and the vector-space
operations are pointwise.
\end{proposition}
\begin{proof}
$\BF$ is a real Banach space (\Cref{prop:banach-space}), so addition and real scaling are
continuous; condensation preserves finite products (\Cref{thm:full-faithful}), so it carries
these to morphisms $\cond{\BF}\times\cond{\BF}\to\cond{\BF}$ and
$\cond{\R}\times\cond{\BF}\to\cond{\BF}$ satisfying the vector-space axioms internally. The
description of $\cond{\BF}(S)$ is \Cref{def:condensation}; $\Cont(S,\BF)$ with $S$ compact and
$\BF$ Banach is complete in the sup-norm, hence itself Banach.
\end{proof}

We now assemble the dynamics. Let $\Aut(\qA)$ carry the topology of strong convergence: a net
$\beta_i\to\beta$ iff $\beta_i(A)\to\beta(A)$ in norm for every $A\in\qA$. Because $\qA$ is
separable, this topology is metrizable and $\Aut(\qA)$ is a Polish group, in particular
compactly generated, so \Cref{thm:full-faithful} applies to it.

\begin{theorem}[Condensed dynamics morphism, I-B]
\label{thm:condensed-dynamics}
Let $\Ffun$ be an $\Ffun$-function, $\qA$ the associated quasi-local algebra, and
$\Aut(\qA)$ its automorphism group in the strong topology.
\begin{enumerate}
\item $\cond{\BF}$ is a condensed $\R$-vector space (\Cref{prop:condensed-BF}).
\item The assignment $D:(t,\Phi)\mapsto\tau^\Phi_t$ is a continuous map
  $\R\times\BF\to\Aut(\qA)$.
\item Consequently $D$ condenses to a morphism of condensed sets
  \[
    \cond{D}\;:\;\cond{\R}\times\cond{\BF}\;\longrightarrow\;\cond{\Aut(\qA)},
    \qquad \cond{\R}\times\cond{\BF}\cong\cond{\R\times\BF},
  \]
  and $\cond{D}$ is a morphism of light condensed sets whenever the interactions are drawn
  from a separable closed subspace of $\BF$ (for instance the finite-range interactions, or
  those with a fixed countable generating family of clusters).
\end{enumerate}
\end{theorem}
\begin{proof}
Part~(1) is \Cref{prop:condensed-BF}. For part~(2), fix $A\in\qA$ and $\varepsilon>0$; we show
$(t,\Phi)\mapsto\tau^\Phi_t(A)$ is continuous, which is the definition of strong continuity of
$D$. Choose a local $A'$ with $\norm{A-A'}<\varepsilon/3$ (possible since $\qA_\loc$ is
dense). Automorphisms are isometric, so
$\norm{\tau^\Phi_t(A)-\tau^{\Phi'}_{t'}(A)}\le
2\norm{A-A'}+\norm{\tau^\Phi_t(A')-\tau^{\Phi'}_{t'}(A')}<\tfrac{2\varepsilon}{3}
+\norm{\tau^\Phi_t(A')-\tau^{\Phi'}_{t'}(A')}$.
By \Cref{lem:dynamics-lipschitz} the second term is small when $\lvert t-t'\rvert$ and
$\Fnorm{\Phi-\Phi'}$ are small, on any ball $\Fnorm{\Phi},\Fnorm{\Phi'}\le B$. Hence $D$ is
continuous at $(t,\Phi)$. Part~(3): $\R\times\BF$ is metrizable, hence compactly generated,
and so is $\Aut(\qA)$ by separability of $\qA$; \Cref{thm:full-faithful} turns the continuous
$D$ into a morphism $\cond{D}$ and identifies $\cond{\R\times\BF}$ with
$\cond{\R}\times\cond{\BF}$. When the interactions lie in a separable closed subspace $V_0
\subseteq\BF$, both $\R\times V_0$ and $\Aut(\qA)$ are separable metrizable, so the restricted
morphism is one of light condensed sets~\cite{clausen-scholze-analytic-stacks}.
\end{proof}

\Cref{thm:condensed-dynamics} is the sense in which ``the dynamics is a morphism, not just an
assignment.'' The diagram
\[
\begin{tikzcd}[column sep=large, row sep=large]
  \cond{\R}\times\cond{\BF} \arrow[r, "\cond{D}"] \arrow[d, "\mathrm{pr}_2"']
    & \cond{\Aut(\qA)} \\
  \cond{\BF} \arrow[ur, "\Phi\ \mapsto\ \cond{\tau^{\Phi}_{\bullet}}"']
\end{tikzcd}
\]
commutes internally in $\Set$: freezing a profinite family of interactions and letting time
vary is a $\cond{\R}$-action on $\cond{\Aut(\qA)}$ over $\cond{\BF}$. This action is the object
that Part~IV group-completes on the invertible gapped sector.

\begin{remark}[Why the Banach topology is not negotiable]
\label{rem:banach-topology}
The continuity in \Cref{thm:condensed-dynamics}(2) is stated for the $\Fnorm{\cdot}$-topology
on $\BF$, and \Cref{lem:dynamics-lipschitz} genuinely uses that norm: the Lieb--Robinson
factor $e^{2BC_\Ffun\lvert t\rvert}$ and the weight $\norm{\Ffun}$ are finite only because
$\Fnorm{\Phi-\Psi}$ measures interaction differences against the same $\Ffun$-function. A
weaker topology (pointwise convergence of $\Phi(X)$ for each cluster $X$, say) does not
control the tail sums and does not make $D$ continuous. This is why the ground floor of the
program is $\cond{\BF}$ for a fixed $\Ffun$ and not the condensation of some coarser space of
formal interactions; long-range interactions outside every $\Ffun$-function class can violate
the finite-velocity bound and fall outside the theorem.
\end{remark}

\section{Uniform Lieb--Robinson bounds over a profinite base}
\label{sec:uniform}

The previous section made the dynamics a morphism out of $\cond{\BF}$. We now show that when
one plugs in an actual profinite family, the Lieb--Robinson estimate holds simultaneously and
uniformly, and we identify what an $S$-continuous family of interactions is in concrete terms.

\subsection{Uniformity over the base}
\label{sec:uniform-bound}

\begin{theorem}[Profinite families and the uniform light cone, I-C]
\label{thm:uniform-lr}
Let $S=\varprojlim_i S_i$ be a profinite set and $\Phi_\bullet\in\Cont(S,\BF)=\cond{\BF}(S)$ a
continuous family of interactions.
\begin{enumerate}
\item \emph{(Compatible finite data.)} $\Phi_\bullet$ is a norm limit, in the sup-norm
  $\norm{\cdot}_{\infty,\Ffun}$ of \Cref{prop:condensed-BF}, of locally constant families:
  there are indices $i_1\le i_2\le\cdots$ and interactions
  $\Phi^{(k)}_\bullet$ each factoring through the finite quotient $S\to S_{i_k}$ with
  $\norm{\Phi_\bullet-\Phi^{(k)}_\bullet}_{\infty,\Ffun}\to 0$. Equivalently,
  $\cond{\BF}(S)$ is the completion of $\varinjlim_i\Cont(S_i,\BF)$ in the sup-norm.
\item \emph{(Uniform bound.)} $B\eqdef\sup_{s\in S}\Fnorm{\Phi_s}<\infty$, attained because
  $S$ is compact and $s\mapsto\Fnorm{\Phi_s}$ is continuous.
\item \emph{(Uniform light cone.)} If in addition $\Phi_s\in\mathcal B_{\Ffun_a}$ with
  $B_a\eqdef\sup_{s\in S}\Fanorm{\Phi_s}<\infty$ for the reweighting $\Ffun_a$, then the
  Lieb--Robinson bound of \Cref{prop:velocity} holds for every $s\in S$ simultaneously with
  \emph{one} velocity
  \[
    v \;=\; \frac{2\,B_a\,C_{\Ffun_a}}{a}
  \]
  and prefactors independent of $s$. Consequently the map $s\mapsto\tau^{\Phi_s}_t$ is a
  morphism $\cond{S}\to\cond{\Aut(\qA)}$ obtained by restricting $\cond{D}$ of
  \Cref{thm:condensed-dynamics} along $\Phi_\bullet$, and it factors through the sub-condensed
  set of automorphisms with light cone of slope at most $v$.
\end{enumerate}
\end{theorem}
\begin{proof}
\emph{(1)} A continuous map from a profinite (compact, totally disconnected) space to a metric
space is a uniform limit of locally constant maps. Concretely, fix $k$ and use uniform
continuity of $\Phi_\bullet$ on the compact $S$: there is a finite clopen partition of $S$ on
which $\Phi_\bullet$ varies by at most $1/k$ in $\Fnorm{\cdot}$; refining, the partition is
pulled back from some finite quotient $S_{i_k}$, and choosing one value of $\Phi_\bullet$ per
block defines $\Phi^{(k)}_\bullet$ factoring through $S_{i_k}$ with
$\norm{\Phi_\bullet-\Phi^{(k)}_\bullet}_{\infty,\Ffun}\le 1/k$. Locally constant families are
exactly the elements of $\varinjlim_i\Cont(S_i,\BF)$ (a locally constant map factors through a
finite quotient), and their sup-norm closure is all of $\Cont(S,\BF)$ by the density just
shown; completeness of $\Cont(S,\BF)$ (\Cref{prop:condensed-BF}) identifies it with the
completion of the colimit.

\emph{(2)} $s\mapsto\Fnorm{\Phi_s}$ is continuous because $\Phi_\bullet$ is continuous and the
norm is $1$-Lipschitz; a continuous real function on a compact set is bounded and attains its
supremum.

\emph{(3)} The estimate of \Cref{prop:velocity} depends on the interaction \emph{only} through
$\Fanorm{\Phi}$, entering the velocity $v=2\Fanorm{\Phi}C_{\Ffun_a}/a$ and the prefactor not at
all. Replacing $\Fanorm{\Phi_s}$ by its uniform bound $B_a$ gives a single velocity
$v=2B_aC_{\Ffun_a}/a$ valid for all $s$; the prefactor $2\norm{A}\norm{B}\norm{\Ffun}
C_{\Ffun_a}^{-1}\min\{\lvert X\rvert,\lvert Y\rvert\}$ is already $s$-independent. Thus every
$\tau^{\Phi_s}_t$ obeys the same light-cone estimate. That $s\mapsto\tau^{\Phi_s}_t$ is a
morphism of condensed sets is \Cref{thm:condensed-dynamics}(3) precomposed with the point
$\cond{S}\to\cond{\BF}$ named by $\Phi_\bullet$; the light-cone slope bound is preserved
because it holds pointwise with the uniform $v$.
\end{proof}

The three parts say, in order: an $S$-continuous family is a compatible tower of
finite-resolution families; such a family is automatically norm-bounded; and a norm bound is
exactly what makes the whole family share one light cone. The last point is the reason the
condensed reformulation is not merely cosmetic. In the topological category one would have a
map $S\to\Int_{d,G}$ and a velocity function $s\mapsto v(s)$ with no guarantee of a common
bound; the Banach structure of $\BF$ upgrades ``finite velocity at each $s$'' to ``one finite
velocity for all $s$'', which is what a morphism into a condensed set of quasi-local dynamics
requires.

The uniformity is compatible with change of probe, which is what makes it useful for descent.

\begin{proposition}[Naturality of the dynamics over the base]
\label{prop:naturality}
For a continuous map $\varphi:S'\to S$ of profinite sets, pullback of families
$\varphi^*:\cond{\BF}(S)\to\cond{\BF}(S')$, $\Phi_\bullet\mapsto\Phi_{\varphi(\,\cdot\,)}$,
commutes with the passage to dynamics: for each fixed $t$ the square
\[
\begin{tikzcd}[column sep=huge, row sep=large]
  \cond{\BF}(S) \arrow[r, "\tau_t"] \arrow[d, "\varphi^*"'] & \cond{\Aut(\qA)}(S) \arrow[d, "\varphi^*"] \\
  \cond{\BF}(S') \arrow[r, "\tau_t"'] & \cond{\Aut(\qA)}(S')
\end{tikzcd}
\]
commutes, and the uniform bound transfers, $\sup_{s'\in S'}\Fnorm{\Phi_{\varphi(s')}}
\le\sup_{s\in S}\Fnorm{\Phi_s}$, so the pulled-back family has light cone of slope at most that
of the original. In particular the finite-quotient approximations of \Cref{thm:uniform-lr}(1)
form a compatible system of quasi-local dynamics with a common velocity.
\end{proposition}
\begin{proof}
Commutativity is naturality of the morphism $\tau_t$ of condensed sets
(\Cref{thm:condensed-dynamics}) evaluated on $\varphi$: a morphism of condensed sets is a
natural transformation of the underlying functors, and its naturality square on the map
$\varphi:S'\to S$ is exactly the diagram. The bound transfers because
$\{\Phi_{\varphi(s')}:s'\in S'\}\subseteq\{\Phi_s:s\in S\}$ as interactions, so its supremum
$\Ffun$-norm cannot exceed that of the original family; the light-cone slope $v=2B_aC_{\Ffun_a}
/a$ is monotone in the uniform bound $B_a$ by \Cref{thm:uniform-lr}(3).
\end{proof}

\begin{corollary}[Disorder families]
\label{cor:disorder}
Let $\dis=Q^{\Z^d}$ be a profinite disorder hull for a finite local-configuration alphabet
$Q$ (a compact, totally disconnected space); here $Q$ denotes the configuration alphabet, and
$\Ffun$ is reserved for the $\Ffun$-function.
A disordered interaction $\omega\mapsto\Phi_\omega$ with
$\sup_{\omega\in\dis}\Fanorm{\Phi_\omega}<\infty$ is an element $\Phi_\bullet\in\cond{\BF}(\dis)$
and hence, by \Cref{thm:uniform-lr}, has a $\dis$-uniform light cone. In particular the
disorder-averaged and almost-sure Lieb--Robinson velocities coincide with the deterministic
bound $v=2\bigl(\sup_\omega\Fanorm{\Phi_\omega}\bigr)C_{\Ffun_a}/a$.
\end{corollary}
\begin{proof}
The map $\omega\mapsto\Phi_\omega$ is continuous into $\BF$ by hypothesis (uniform-norm
continuity), so it is an element of $\Cont(\dis,\BF)=\cond{\BF}(\dis)$; apply
\Cref{thm:uniform-lr}(3) with $S=\dis$. The averaged and almost-sure velocities are bounded by
the uniform one because the estimate is deterministic and holds for every $\omega$.
\end{proof}

This corollary is the locality-layer appearance of the disorder thread that runs through the
series: the profinite probe $\dis$ \emph{is} the physical configuration space, not a
repackaging of a smooth manifold, and \Cref{cor:disorder} is what \Cref{conj:descent} below
proposes to sharpen into a descent statement.

\subsection{Uniform clustering: the bridge to the gapped substack}
\label{sec:uniform-clustering}

Part~III needs not just the light cone but its standard spectral consequence: a uniform gap
forces uniform exponential decay of ground-state correlations. We state the version that our
uniformity provides; the single-Hamiltonian statement is due to
Hastings--Koma~\cite{hastings-koma} and Nachtergaele--Sims~\cite{nachtergaele-sims-clustering}.

\begin{proposition}[Uniform exponential clustering, input to Part III]
\label{prop:clustering}
Let $\Phi_\bullet\in\cond{\mathcal B_{\Ffun_a}}(S)$ over a profinite $S$ with uniform bound
$B_a=\sup_s\Fanorm{\Phi_s}<\infty$, and suppose the family is uniformly gapped: there is
$\gapbound>0$ with $\gapof{H^{\Phi_s}}\ge\gapbound$ for all $s\in S$, where $\gapof{\cdot}$ is
the spectral gap above the ground state of the infinite-volume GNS Hamiltonian. Then there are
constants $C,\mu>0$, depending only on $B_a$, $C_{\Ffun_a}$, $a$, and $\gapbound$---and in
particular not on $s$---such that for all local $A,B$ with disjoint supports and all $s\in S$,
\[
  \bigl\lvert\langle\Omega_s\lvert AB\rvert\Omega_s\rangle
    -\langle\Omega_s\lvert A\rvert\Omega_s\rangle
     \langle\Omega_s\lvert B\rvert\Omega_s\rangle\bigr\rvert
  \;\le\; C\,\norm{A}\,\norm{B}\,e^{-\mu\,\dist(\supp A,\supp B)},
\]
where $\Omega_s$ is the ground state of $H^{\Phi_s}$.
\end{proposition}
\begin{proof}
For a fixed $s$ this is the Hastings--Koma / Nachtergaele--Sims exponential clustering
theorem~\cite{hastings-koma,nachtergaele-sims-clustering}: a spectral gap $\gapbound$ together
with a Lieb--Robinson velocity $v$ yields correlation decay with rate
$\mu=\mathrm{const}\cdot\gapbound/(v+\gapbound)$ and constant $C$ depending on the same data.
The only point to check is uniformity. By \Cref{thm:uniform-lr}(3) the velocity $v=2B_a
C_{\Ffun_a}/a$ and the Lieb--Robinson prefactors are the same for every $s$; by hypothesis the
gap lower bound $\gapbound$ is the same for every $s$. The clustering rate and constant are
explicit functions of $(v,\gapbound,C_{\Ffun_a},a)$ only, so they are $s$-independent.
\end{proof}

\Cref{prop:clustering} is stated conditionally on a uniform gap because deciding the gap is not
a locality question and, by the Cubitt--P\'erez-Garc\'ia--Wolf undecidability
theorem~\cite{cpw-undecidable}\footnote{Cited for scope, not used: there is no algorithm
deciding whether a translation-invariant nearest-neighbour family is gapped. Part~III therefore
treats ``uniformly gapped'' as the hypothesis defining $\Gap_{d,G}$, never as a conclusion.},
cannot be. Our contribution is exactly the clause ``and in particular not on $s$'': the
locality estimates are uniform, so whatever gap hypothesis Part~III imposes, its clustering
consequence is automatically uniform over the base.

\section{Quasi-adiabatic continuation and the class \texorpdfstring{$\Weq$}{W}}
\label{sec:qac}

The moduli-stack picture inverts a class $\Weq$ of gapped adiabatic, quasi-local equivalences
to define phases. At the level of a single gapped path the rigorous content of $\Weq$ is
quasi-adiabatic continuation (Hastings--Wen~\cite{hastings-wen-qac}) and the automorphic
equivalence of Bachmann--Michalakis--Nachtergaele--Sims~\cite{bmns-automorphic}. We recall
these as the morphism-level tools they are, staying within what is proved. Throughout we keep
the Heisenberg picture fixed in \Cref{sec:lr-dynamics}, where automorphisms act by
$A\mapsto U^*AU$, so that $\tau^\Phi_t(A)=e^{itH^\Phi}Ae^{-itH^\Phi}$; this convention is what
puts the positive sign in the commutator flows below.

\begin{definition}[Quasi-adiabatic generator]
\label{def:qac}
Let $[0,1]\ni r\mapsto\Phi(r)\in\mathcal B_{\Ffun_a}$ be a norm-$C^1$ path of interactions,
uniformly gapped with gap $\ge\gapbound>0$, and let $P_r$ be the ground-state projection of
$H^{\Phi(r)}$. The \emph{quasi-adiabatic generator} is
\[
  \mathcal D(r)\;=\;\int_{-\infty}^{\infty} W_\gamma(t)\;
     \tau^{\Phi(r)}_t\!\Bigl(\tfrac{d}{dr}H^{\Phi(r)}\Bigr)\,dt,
\]
where $W_\gamma$ is a Hastings weight function whose Fourier transform is supported outside
$(-\gapbound,\gapbound)$ and decays faster than any polynomial (indeed almost exponentially).
\end{definition}

\begin{proposition}[Quasi-adiabatic continuation generates $\Weq$~\cite{hastings-wen-qac,bmns-automorphic}]
\label{prop:qac}
With the hypotheses of \Cref{def:qac}:
\begin{enumerate}
\item $\mathcal D(r)$ is a quasi-local self-adjoint operator: its finite-volume
  approximations obey a Lieb--Robinson bound for a reweighted $\Ffun$-function
  $\Ffun_{a'}$ with $a'<a$, and $\norm{\mathcal D(r)}_{\Ffun_{a'}}$ is bounded uniformly in
  $r$ in terms of $\gapbound$ and $\sup_r\Fanorm{\tfrac{d}{dr}\Phi(r)}$.
\item The flow $\alpha_r$ generated by $\{\mathcal D(r)\}$ (the solution of
  $\tfrac{d}{dr}\alpha_r=i\,\alpha_r\circ[\mathcal D(r),\,\cdot\,]$, $\alpha_0=\id$) is a
  strongly continuous family of $*$-automorphisms of $\qA$ with $\alpha_r(P_0)=P_r$ for all
  $r$; equivalently $\alpha_r$ intertwines the ground states along the path.
\item Two interactions connected by such a uniformly gapped path are
  \emph{automorphically equivalent}: there is a quasi-local automorphism $\alpha_1$, obeying a
  Lieb--Robinson bound, with $\alpha_1$ mapping the ground state of $H^{\Phi(0)}$ to that of
  $H^{\Phi(1)}$. This automorphism is the generator of $\Weq$ for the pair.
\end{enumerate}
\end{proposition}
\begin{proof}[Proof (citation)]
Parts~(1)--(2) are the construction of Hastings--Wen~\cite{hastings-wen-qac}, made quantitative
in the $\Ffun$-function calculus of Nachtergaele--Sims--Young~\cite{nsy-quasilocality-1}: the
quasi-locality of $\mathcal D(r)$ follows from the Lieb--Robinson bound of \Cref{thm:banach}
for $\tau^{\Phi(r)}_t$ and the fast decay of $W_\gamma$, and the intertwining property is the
defining feature of the Hastings weight. Part~(3) is the automorphic equivalence theorem of
Bachmann--Michalakis--Nachtergaele--Sims~\cite{bmns-automorphic}: integrating the flow of
part~(2) over $r\in[0,1]$ yields the quasi-local $\alpha_1$. We claim nothing beyond these
results; in particular we do not assert that $\alpha_1$ assembles over a profinite base into a
condensed automorphism, which is the content of \Cref{conj:group}.
\end{proof}

\begin{remark}[What $\Weq$ is, and is not, here]
\label{rem:W}
$\Weq$ is part of the \emph{definition} of the phase problem, not a theorem: isomorphism,
gapped homotopy, stable equivalence, and $K$-theoretic equivalence need not coincide, and the
program must say which one it inverts. \Cref{prop:qac} pins down the adiabatic, quasi-local
generators of $\Weq$ for a single gapped path. Turning these generators into an honest
morphism structure over $\cond{\BF}$ (a condensed group acting on $\Ham_{d,G}$ with
quasi-adiabatic continuation as internal path-lifting) is beyond the present estimates and is
posed as \Cref{conj:group}.
\end{remark}

\section{Conjectures: toward the condensed stack}
\label{sec:conjectures}

The theorems above make $\cond{\BF}$ a condensed $\R$-vector space and the dynamics a
morphism out of it, uniformly over profinite probes. They do \emph{not} establish that the
whole moduli problem $\Ham_{d,G}$ (which remembers gauge and quasi-local automorphisms, not
just interactions) is a condensed stack. We isolate the three statements that the
locality layer would need, and label them as conjectures with stable numbers.

\begin{conjecture}[Condensed higher stack, I-1]
\label{conj:stack}
The moduli problem $\Ham_{d,G}$ of $G$-symmetric quasi-local Hamiltonians is a condensed
higher stack, not merely a condensed set: the assignment sending a profinite probe $S$ to the
groupoid (or $\infty$-groupoid) of $S$-families of Hamiltonians and their gauge and
quasi-local automorphisms satisfies descent along finite jointly surjective covers of profinite
sets. Its $\pi_0$ over the point recovers the condensed set $\cond{\BF}$ modulo relabelling,
and its automorphism sheaf is generated by the quasi-local flows of \Cref{prop:qac}.
\end{conjecture}

What is proved here is the underlying-condensed-set statement (\Cref{prop:condensed-BF}) and
the morphism property of dynamics (\Cref{thm:condensed-dynamics}); the gluing of the
automorphism groupoid along covers is open. We flag, as the knowledge base insists, that this
stack structure must not be asserted as a theorem anywhere in the series.

\begin{conjecture}[Condensed group of quasi-local automorphisms, I-2]
\label{conj:group}
The quasi-local automorphisms generated by time-dependent interactions in a fixed
$\mathcal B_{\Ffun_a}$-ball, with the Lieb--Robinson bound of \Cref{prop:qac}(1), assemble
into a condensed group $\underline{\mathcal G}_{\Ffun_a}$ acting on $\cond{\BF}$ (and on
$\Ham_{d,G}$). Under this action, quasi-adiabatic continuation is an internal path-lifting:
a uniformly gapped $S$-family of paths lifts to an $S$-family of automorphisms in
$\underline{\mathcal G}_{\Ffun_a}$, continuously in the probe.
\end{conjecture}

The obstruction to a proof is exactly the uniformity of the quasi-adiabatic construction over
a base: \Cref{prop:qac} produces $\alpha_1$ for each path, and \Cref{thm:uniform-lr} makes the
underlying Lieb--Robinson data uniform, but assembling the $\alpha_1$ into a continuous
$S$-family of automorphisms with a common quasi-locality modulus is not contained in the cited
theorems.

\begin{conjecture}[Descent of Lieb--Robinson estimates along a disorder hull, I-3]
\label{conj:descent}
Let $\dis=Q^{\Z^d}$ be a profinite disorder hull with its $\Z^d$-action, and
$\Phi_\bullet\in\cond{\BF}(\dis)$ a disordered family. The Lieb--Robinson estimates of
\Cref{thm:uniform-lr} satisfy descent along the finite quotients $\dis\to\dis_i$: the
quasi-local dynamics over $\dis$ is the limit of the quasi-local dynamics over the finite
quotients in a way that determines all quasi-local invariants from finite-resolution data.
Equivalently, the sheaf $S\mapsto\{\text{quasi-local dynamics over }S\}$ is the right Kan
extension of its restriction to finite quotients of $\dis$.
\end{conjecture}

\Cref{cor:disorder} is the unconditional shadow of \Cref{conj:descent}: it gives the uniform
velocity over $\dis$, which is the $\pi_0$-level statement. The descent claim is the assertion
that the entire quasi-local structure, not just the velocity, is finite-resolution data; this
is what Part~II would use to make the crossed-product observable algebra $C(\dis)\rtimes\Z^d$
functorial in the finite quotients.

\section{The light cone in the transverse-field Ising chain}
\label{sec:numerics}

The accompanying Haskell package \texttt{src/lieb-robinson-locality/} makes the light cone of
\Cref{ex:tfim} visible and fits a velocity. We summarize what it computes; the code is the
authoritative specification. It is available at
\href{https://github.com/YonedaAI/topological-phases-of-matter}{github.com/YonedaAI/topological-phases-of-matter}
under \texttt{src/lieb-robinson-locality/}, so the numbers below reproduce from a single
\texttt{ghc} (or \texttt{cabal}) build.

The \texttt{FFunction} module implements $\Ffun$-functions as first-class values with
combinators for power-law decay $\Ffun_\varepsilon(r)=(1+r)^{-(d+\varepsilon)}$, exponential
decay, and the reweighting $\Ffun\mapsto\Ffun_a$ of \Cref{lem:reweight}; it exposes numerical
estimates of $\norm{\Ffun}$ and $C_\Ffun$ on a truncated $\Z$. The \texttt{Ising} module
builds the finite-volume transverse-field Ising chain, evolves a local operator $A=\sigma^x_0$
by the exact matrix exponential $\tau_t(A)=e^{itH}Ae^{-itH}$, and computes the commutator norm
\[
  c(x,t)\;=\;\bigl\lVert[\tau_t(\sigma^x_0),\,\sigma^x_x]\bigr\rVert
\]
on a grid of sites $x$ and times $t$. Plotting the level sets of $c(x,t)$ shows the
characteristic cone $\lvert x\rvert\approx v_{\mathrm{fit}}\,t$ outside of which $c$ is
exponentially small; a least-squares fit through the origin of the front position---the
largest site $x$ whose commutator norm reaches $10\%$ of the row maximum at that time---against
$t$ recovers a group velocity $v_{\mathrm{fit}}$ consistent with the analytic bound of
\Cref{prop:velocity}. For
the isotropic point $J=h=1$ the fitted velocity is of order the known
$v_{\mathrm{LR}}=2\max(J,h)$ scale, and it stays finite and bounded as $(J,h)$ range over a
compact box, illustrating the uniform-over-the-base content of \Cref{thm:uniform-lr}: the
velocity does not blow up as the parameters vary within bounds.

The \texttt{Properties} module states QuickCheck properties that formalize the paper's
elementary claims: monotonicity of $\Ffun$-functions (power-law and exponential), the two
reweighting inequalities $\norm{\Ffun_a}\le\norm{\Ffun}$ and $C_{\Ffun_a}\le C_\Ffun$ of
\Cref{lem:reweight} (both tested on the truncated constants), the subadditivity
$\Fnorm{\Phi+\Psi}\le\Fnorm{\Phi}+\Fnorm{\Psi}$ of the interaction norm, and the positivity and
monotonicity in the norm bound of the analytic velocity of \Cref{prop:velocity}.
\texttt{Main.hs} runs the simulation, prints the fitted velocity alongside the analytic
Lieb--Robinson upper bound, checks deterministically that the former does not exceed the
latter, and exits with status zero.

\begin{remark}
The simulation is finite-volume and therefore does not, and cannot, verify the infinite-volume
statements of \Cref{thm:banach,thm:condensed-dynamics}; it verifies the finite-volume
Lieb--Robinson estimate of \Cref{thm:lr-bound}, which is the uniform-in-$\Lambda$ input those
theorems pass to the limit. That is the honest scope of a numerical check: it exhibits the
mechanism (a finite velocity, stable under parameter variation) whose analytic control is the
content of the theorems.
\end{remark}

\section{Discussion}
\label{sec:discussion}

\paragraph{What the condensed packaging buys.}
Nothing in \Cref{thm:banach} is new as analysis; it is Nachtergaele--Sims--Young in a fixed
notation. The step that is genuinely a step is \Cref{thm:condensed-dynamics} together with
\Cref{thm:uniform-lr}: the Lieb--Robinson estimate is not merely a bound but a
\emph{continuity} statement in the Banach norm, and continuity is exactly what condensation
requires. Once one sees this, the dynamics is a morphism $\cond{\R}\times
\cond{\BF}\to\cond{\Aut(\qA)}$ for free, and profinite families (finite-volume towers,
disorder hulls, inverse limits of parameter spaces) inherit a uniform light cone with no extra
work. The payoff is organizational, as the program concedes throughout: we do not compute a new
invariant, we make the invariants of later parts live over a base that includes disorder and
continuous families on the same footing as single Hamiltonians.

\paragraph{Limitations.}
Three boundaries are worth stating plainly. First, everything is tied to a fixed
$\Ffun$-function; interactions with slower-than-$\Ffun$ tails, and genuinely long-range models,
are outside the theorems, and for some of them the finite-velocity picture is known to
fail. Second, the condensed \emph{set} structure is all that is proved: the higher-stack
structure of $\Ham_{d,G}$ (\Cref{conj:stack}) and the condensed group of automorphisms
(\Cref{conj:group}) are conjectural, and we have been careful never to use them as if proved.
Third, the light-condensed reduction requires separability; for a non-separable interaction
space one is in the full condensed formalism, with its attendant set-theoretic care, and the
statements about $\Aut(\qA)$ being Polish use separability of $\qA$.

\paragraph{Where the thread continues.}
The automorphic equivalence of \Cref{prop:qac} is the generator of the class $\Weq$ that
Part~III inverts to form the phase groupoid; the uniform clustering of \Cref{prop:clustering}
is the correlation-decay input to the stability theory there; and the condensed morphism of
\Cref{thm:condensed-dynamics} is what lets Part~IV's parametrized families be morphisms of
condensed spectra. The disorder corollary (\Cref{cor:disorder}) and its conjectural
sharpening (\Cref{conj:descent}) feed the crossed-product picture of Part~II. In each case the
locality layer supplies uniformity, and uniformity is what makes the later constructions
functorial in the probe rather than merely pointwise.

\section{Conclusion}
\label{sec:conclusion}

We have built the analytic ground floor of the condensed-mathematics program for topological
phases. The $\Ffun$-normed $G$-symmetric interactions form a real Banach space $\BF$
(\Cref{thm:banach}), its condensation $\cond{\BF}$ is a condensed $\R$-vector space, and the
Heisenberg dynamics is a genuine morphism $\cond{\R}\times\cond{\BF}\to\cond{\Aut(\qA)}$ of
condensed sets (\Cref{thm:condensed-dynamics}) because the Lieb--Robinson bound is a continuity
statement in the Banach norm. Over a profinite base an $S$-continuous family of interactions is
a compatible tower of finite-resolution families, and a uniform $\Ffun$-norm bound forces a
single Lieb--Robinson velocity and one light cone across the whole base
(\Cref{thm:uniform-lr}), with uniform exponential clustering as a corollary under a uniform gap
(\Cref{prop:clustering}) and quasi-adiabatic continuation supplying the generators of the
equivalence class $\Weq$ (\Cref{prop:qac}). The three conjectures I-1, I-2, I-3 mark exactly
where today's estimates stop: the higher-stack structure, the condensed group of automorphisms,
and descent along a disorder hull are the locality-layer instances of the program's Master
Conjecture, and they are stated as conjectures precisely because the cited technology does not
yet reach them. The subsequent parts build positivity, the gap, effective field theories, and
realizability on top of this floor; each of them takes ``$S$-continuous family with a uniform
light cone'' as a primitive, and that primitive is what this paper makes precise.

\appendix
\section{Constants in the Lieb--Robinson bound}
\label{app:constants}

For reference we collect the constants used above, all for a fixed $\Ffun$-function $\Ffun$ on
$(L,\dist)$ and its reweighting $\Ffun_a$.
\begin{itemize}
\item $\norm{\Ffun}$ (uniform summability) and $C_\Ffun$ (convolution) are the two constants
  of \Cref{def:Ffun}; reweighting gives $\norm{\Ffun_a}\le\norm{\Ffun}$ and $C_{\Ffun_a}\le
  C_\Ffun$ (\Cref{lem:reweight}).
\item The commutator prefactor in \Cref{thm:lr-bound} is $2\norm{A}\norm{B}/C_\Ffun$ and the
  time factor is $e^{2\Fnorm{\Phi}C_\Ffun\lvert t\rvert}-1$.
\item The velocity in \Cref{prop:velocity} is $v=2\Fanorm{\Phi}C_{\Ffun_a}/a$; minimizing the
  bound over $a>0$ (subject to $\Phi\in\mathcal B_{\Ffun_a}$) gives the sharpest cone the
  method provides.
\item The Lipschitz constant of the dynamics in the interaction
  (\Cref{lem:dynamics-lipschitz}) is $K=2\norm{A}\,\lvert\supp A\rvert\,\norm{\Ffun}\,T\,
  e^{2BC_\Ffun T}$ on the ball $\Fnorm{\cdot}\le B$ and time horizon $T$.
\item The uniform velocity over a profinite base (\Cref{thm:uniform-lr}) is
  $v=2B_aC_{\Ffun_a}/a$ with $B_a=\sup_{s\in S}\Fanorm{\Phi_s}$.
\end{itemize}

\addcontentsline{toc}{section}{References}
\bibliographystyle{plain}
\bibliography{references}

\end{document}
